\begin{proposition}\label{prop:shauder-theorem}
    Если $X$ -- нормированное, $Y$ -- банахово и $T \in \LL(X, Y)$, то следующие условия эквивалентны:
    \begin{itemize}
        \item $T$ -- компактный;
        \item $T^*$ -- компактный.
    \end{itemize}
\end{proposition}
\begin{proof}
    Пусть $T$ -- компактен.
    Рассмотрим $T^* B_1^*(0)$ и выберем в нём произвольную последовательность $\{f_n\} \subset X^*$.
    То есть $f_n = T^* g_n$, где $g_n \in B_1^*(0) \subset Y^*$.
    Для вполне ограниченности $T^* B_1^*(0)$ достаточно показать, что мы можем выделить фундаментальную подпоследовательность.
    По определению сопряжённого оператора:
    \[
        f_n(x) = (T^* g_n)(x) = g_n(Tx).
    \]
    Тогда
    \[
        \|f_n - f_m\| = \sup_{\|x\| \leq 1} |f_n(x) - f_m(x)| = \sup_{\|x\| \leq 1}|g_n(Tx) - g_m(Tx)|
    \]
    В силу компактности $T$ множество $TB_1(0)$ -- вполне ограничено, а значит $K = [TB_1(0)]$ -- компакт в $Y$.
    Рассмотрим $C(K)$.
    Заметим, что всякий $g_n \in C(K)$, и
    \[
        \|g_n\|_{C(K)} = \max_{y \in K}|g_n(y)| = \sup_{x \in B_1(0)} |g_n(Tx)| = \sup_{x \in B_1(0)}|f_n(x)| = \|f_n\|.
    \]
    А значит
    \[
        \|f_n - f_m\| = \|g_n - g_m\|_{C(K)}.
    \]
    Пусть $S = \{g_n\} \subset C(K)$.
    Проверим условия критерия Арцела-Асколи.
    Очевидно, что $S$ -- ограничено (поскольку $\|g_n\|_{C(K)} = \|f_n\| \leq \|T\|$).
    Так как
    \[
        |g_n(y) - g_n(z)| = |g_n(y - z)| \leq \|g_n\| \|y - z\|,
    \]
    и, в силу того, что $g_n \in B_1^*(0)$ выполняется $\|g_n\| \leq 1$, то
    \[
        |g_n(y) - g_n(z)| \leq \|y - z\|.
    \]
    Что и даёт нам равностепенную непрерывность, а значит и вполне ограниченность $S$ -- а значит и вполне ограниченность $\{f_n\}$ -- из чего, в силу произвольности её выбора, следует вполне ограниченность $T^* B_1^*(0)$ -- то есть компактность оператора $T^*$.

    Пусть $T^*$ -- компактный оператор.
    По только что доказанному $T^{**} \in \LL(X^{**}, Y^{**})$ является компактным оператором.
    Поскольку $T = \Phi_Y^{-1} \circ T^{**} \circ \Phi_X$ (где $\Phi_X, \Phi_Y$ -- канонические изометрии Банаха), то он компактен ($TB_1(0) = \Phi_Y^{-1}(T^{**}\Phi_X(B_1(0)))$ и изометричные операторы сохраняют вполне ограниченность (и просто ограниченность)).
\end{proof}
\begin{example}[Уравнение Фредгольма и критерий Фредгольма]
     Пусть $\lambda \in \Cm \setminus \{0\}$, $X$ -- банахово пространство и $A \in \LL(X)$.
     Рассмотрим уравнение (так называемое уравнение Фредгольма):
     \[
        Ax = \lambda x + y \Longleftrightarrow A_\lambda x = y.
     \]
\end{example}
\begin{example}
    Пусть $T\colon \ell^1 \to \ell^1$ -- оператор левого сдвига, то есть $Tx(n) = x(n + 1)$.
    Заметим, что $\im T = \ell^1$, но при этом $\ker T = \Lin\{e_1\}$ и $\lambda = 0$ -- собственное значение оператора $T$. Таким образом, для произвольных непрерывных операторов не выполняются матричные соотношения, связывающие размерность ядра и образа. Тем не менее, для компактных операторов это верно (в некотором смысле).
\end{example}
\begin{theorem}
    Пусть $X$ -- банахово пространство, $A \in \LL(X)$ -- компактный оператор и $\lambda \neq 0$ -- собственное число оператора $A$ (то есть $\ker A_\lambda \neq 0$). Тогда $\im A_\lambda \neq X$.
\end{theorem}
\begin{proof}
    Пусть $\im A_\lambda = X$.
    Тогда существует $x_1 \in \ker A_\lambda \setminus \{0\}$, существует $x_2 \in X\colon x_1 = A_\lambda x_2 \neq 0$, и так далее (то есть мы выбираем $x_3\colon \ A_\lambda x_3 = x_2\ldots$).
    Таким образом мы построили последовательность $\{x_m\}$ такую, что $x_m \in \ker A_\lambda^m$, но $x_m \notin \ker A_\lambda^{m - 1}$.
    Пусть $L_m = \ker A_\lambda^m$. Мы получаем строго возрастающую последовательность подпространств $L_m$.

    По лемме Рисса о почти перпендикуляре $\forall m \geq 2$ существует $y_m \in L_m\colon \|y_m\| = 1$ такой, что $\rho(y_m, L_{m - 1}) \geq \frac{1}{2}$.

    Поскольку $A$ -- компактен, то $\{Ay_m\}$ -- вполне ограничено и, следовательно, имеет фундаментальную подпоследовательность.
    Но, при этом, если $m > n \geq 2$, то
    \[
        \|Ay_m - Ay_n\| = \|A_\lambda y_m - A_\lambda y_n - \lambda y_n + \lambda y_m\| = |\lambda| \left\|y_m - y_n + \frac{A_\lambda y_m - A_\lambda y_n}{\lambda}\right\|.
    \]
    Заметим, что $A_\lambda y_m \in L_{m - 1}$.
    То есть $\frac{A_\lambda y_m - A_\lambda y_n}{\lambda} - y_n \in L_{m - 1} + L_{n - 1} + L_n \subset L_{m - 1}$.
    Тогда, в силу выбора $y_m$, мы получаем
    \[
        \|Ay_m - Ay_n\| \geq \frac{|\lambda|}{2}.
    \]
    Следовательно, $\{Ay_m\}$ не может иметь фундаментальной подпоследовательности -- противоречие.
\end{proof}
\begin{note}
    В частности, заметим, что если для компактного $A$ $\ker A_\lambda \neq 0$ и $\im A_\lambda\neq X$, то $\ker A_\lambda^* = \im A_\lambda^\perp \neq 0$.
\end{note}